\chapter{Methodology}
\label{ch:methodology}

The methodology of this research is organized into three phases: (1)~development of an analytical framework from the literature to evaluate existing document protection methods, (2)~design-based research to develop an artifact meeting the design criteria established in Phase~1, and (3)~a within-subjects user study to measure the usability and user acceptance of the developed system.

First, we perform a systemization of knowledge for existing document protection methods by creating an analytical framework similar to those used to evaluate password alternatives~\cite{bonneau_quest_2012} and secure email~\cite{borisov_sok_2021}. The categories and evaluation criteria we present in this work are consistent with how usability-security-deployability frameworks have been developed and applied in prior literature. We apply this framework to compare nine document protection methods across 15 design properties, identifying each method's benefits and limitations. Development of the document encryption framework involved review of related standardized evaluation frameworks, an analysis of technical specifications for each product, and expert consultation to validate decisions.

Second, from the analysis of existing methods, we derive design requirements for a system that is not password-based, not proprietary, does not require recipients to install special software, and encrypts documents individually. We then conduct design-based research to develop a novel passwordless document encryption system leveraging FIDO2 passkeys, the WebAuthn PRF extension, and open-source age encryption to provide end-to-end protection of digital documents without passwords. An iterative development process produced a proof-of-concept implementation, presented as an open-source codebase and web application, which was evaluated against the established design requirements. Benefits include no user memory burden, encryption keys derived from hardware-backed credentials, and no third-party vendor risk.

Third, we evaluate the system through an IRB-approved within-subjects user study (N=21) comparing the passwordless approach to traditional password-based encryption as provided by Adobe Acrobat. Participants completed encryption and sharing tasks with both systems in a lab setting, providing both quantitative and qualitative data. Usability was measured using the System Usability Scale (SUS)~\cite{brooke_sus_1996} and user acceptance using the Van der Laan scale~\cite{vanderlaan_simple_1997}, with the Affinity for Technology Interaction (ATI)~\cite{franke_personal_2019} and Privacy Concern~\cite{langer_information_2018} scales used to identify covariates. Internal consistency of each instrument was assessed using Cronbach's alpha. The Wilcoxon signed-rank test was used for between-condition comparisons due to the within-subjects design and small sample. Participants also answered ten open-ended interview questions covering their experience with each method, perceived advantages and disadvantages, trust and security perceptions, current document-handling practices, willingness to adopt the passwordless method, and suggestions for improvement. The full protocol is provided in Appendix~\ref{ap:instruments}. These semi-structured interviews were analyzed through inductive thematic coding~\cite{braun_thematic_2006} to identify patterns, similarities, and differences in participant experiences across the two conditions. Chapter~\ref{ch:user-study} describes the methodology of the user study in greater detail.
